% ============================================================
% ch1_introduction.tex
% ANSpect O-EMA · MS Thesis · UTD Electrical Engineering
% \input target — no \documentclass or \begin{document}
% STATUS: SKELETON — run 'expand CH-1.X' to flesh out
% v1.0  01-Sep-2026  Split from chapters.tex
% ============================================================

% ============================================================
% CHAPTER 1 — INTRODUCTION
% STATUS: OPEN
% ============================================================

\chapter{Introduction}
\label{ch:introduction}
% TODO:
%   - Brief framing of the research context
%   - Cable broadcast PHY implementation on RFSoC platform
%   - Lead into chapter structure

% ------------------------------------------------------------
\section{Motivation}
\label{sec:motivation}
% TODO:
%   - Growth of FPGA-based SDR systems in cable/broadcast
%   - Gap: no open academic RFSoC J.83B Annex B implementation
%   - Industrial relevance: DOCSIS 3.0 downstream PHY demand
%   - Motivate RTL-first approach over MATLAB/HDL Coder path

% ------------------------------------------------------------
\section{Problem Statement}
\label{sec:problem_statement}
% TODO:
%   - Define the technical problem: full J.83B Annex B TX+RX
%     chain on Zynq UltraScale+ RFSoC 4x2
%   - Identify HDL Coder limitations as a secondary problem
%   - State the scope boundary (what is and is not in this work)

% ------------------------------------------------------------
\section{Contributions}
\label{sec:contributions}
% TODO:
%   - Enumerate thesis contributions as a list
%   Contribution 1: Manual SystemVerilog RTL implementation of
%     full J.83B Annex B chain — motivated by HDL Coder
%     limitations (document decision and outcome)
%   Contribution 2: [NDA-TITLE system] — state generically,
%     no algorithm or implementation detail
%     (e.g., "a complete PYNQ-based loopback testbed for ...")
%   - Cite relevant standards used \cite{NEEDED} % [CITE NEEDED]

% ------------------------------------------------------------
\section{Thesis Organization}
\label{sec:organization}
% TODO:
%   - Chapter-by-chapter roadmap paragraph
%   - CH-2: FPGA/RFSoC background
%   - CH-3: J.83B DSP algorithm background
%   - CH-4: implementation details
%   - CH-5: conclusions and future work
